\scanpage{7}\mspage{3}
\deleted{\unsure{On}} on a $u^2\in V^*$, donc on a une application
\[
\mathrm{Circ}(I)\longrightarrow V^*
\]
de degré $2$. (qui définit donc une structure de \uline{cube}, ayant $V^\times$ comme ens des directions de faces, $\mathrm{Circ}(I)$ comme l'ens des six faces -- en fait ce cube est canoniquement orienté, cf plus bas.
% slanted
\margin{$\circledast$ On peut dire aussi que les six éléments d'ordre $4$ se répartissent par paires $\{\unsure{u},u^{-1}\}$ suivant les trois sous-groupes cycliques d'ordre $4$ \deleted{\unsure{de}} $G=\unsure{\mathfrak{S}|G}_I$, en corr. biunivoque (\unsure{via leurs élts d'ordre 2}) avec \deleted{\unsure{les}} les élts de $V^*$.}
% slanted
\margin{$I$ est l'ens des quatre diagonales du cube,}

L'isom canonique
\[
\underline{\omega}(I)\simeq\underline{\omega}(J)\tag*{$(\unsure{\omega|\alpha|*|**})$}
\]
On a
\[
\mathfrak{S}_{I\,\mathrm{ab}}\xrightarrow{\sim}\mathfrak{S}_{J\,\mathrm{ab}}\simeq\{\pm1\}
\]
les isom. étant donnés par la signature, d'où la possibilité de définir un iso. \deleted{\unsure{can}} \inserted{fonctoriel} $(\omega)$, dont le choix reste soumis cependant à convention (deux choix possibles). Quand on fixe un élément $i\in I$, alors on a
\[
\underline{\omega}_{\unsure{1|i}}(I)\xrightarrow{\sim}\underline{\omega}(\underbrace{I\setminus\{i\}}_{\simeq J})\deleted{\unsure{\times\{i,I\setminus\{i\}\}}}
\]
en associant à tt ordre \inserted{total} de \deleted{$\unsure{J}$} l'ordre \deleted{\unsure{induit}} \inserted{total} de $I$ qui sur $J$ est l'ordre donné, et qui admet $i$ comme \uline{premier} élément (la convention opposée aurait été de prendre $i$ comme \uline{dernier} élément).

\diagram{\begin{tikzcd}[ampersand replacement=\&,row sep=7pt,column sep=5pt,cells={nodes={font=\scriptsize}},execute at end picture={
\draw ([xshift=-2pt,yshift=7pt]opening.north west) .. controls ([yshift=13pt]source.north) and ([yshift=12pt]star.north) .. ([yshift=9pt]cyclic.north east);
\draw ([xshift=-3pt,yshift=5pt]mark.north west)--([xshift=-2pt,yshift=-2pt]norm.south west);
\draw[line width=.8pt] ([xshift=-6pt,yshift=-1pt]mark.north west)--([xshift=-8pt,yshift=-3pt]mark.north west)--([xshift=-7pt,yshift=-4pt]mark.north west);
}]
|[alias=opening]|\ast \&
|[alias=source]|V^*\arrow[r,shift left=2pt,"{x\mapsto\mathrm{Cent}_G(x)}","{\sim}" description] \&
\substack{\text{Sous-groupes}\\\text{cycliques}\\\text{d'ordre }4\\\text{de }G=\mathfrak{S}_I}\arrow[l,shift left=2pt,"{\sim}"'] \&
|[alias=mark,xshift=-1pt,yshift=-8pt]|\unsure{x|\ast}\arrow[d,start anchor=south east,end anchor=north west] \&
|[alias=star]|\substack{V^*\\\deleted{\illegible{1}}}\arrow[rr,leftrightarrow]\arrow[dr,start anchor=south east,end anchor=north west] \& {} \&
|[alias=cyclic]|\substack{\text{Ss-groupes cycliques}\\\text{d'ordre }4\text{ de }G=\mathfrak{S}_I}\arrow[dl,leftrightarrow]\\
\&\substack{\text{l'unique}\\\text{élément }\unsure{\ne|\#}1\\\text{de }\Gamma\cap V}
\& \Gamma\arrow[l,maps to]
\& |[alias=norm,xshift=4pt]|\mathrm{Norm}
\& D\mapsto\substack{\deleted{\unsure{\mathrm{Cent}(D)}}\\\text{élément}\\\text{non nul}\\\text{du centre}\\\text{de }D}\arrow[u,start anchor=north east,end anchor=south west]
\& \substack{\text{Ss-groupes de}\\\text{Sylow }\deleted{\unsure{2}}\text{ de }G=\mathfrak{S}_I\\D)} \&
\end{tikzcd}}
% archive sheet number: 6
